\cotefile{156-8}
\scanpage{58}
\mspage{49}
% archive sheet number: 58
Il y a une difficulté avec $\unsure{f^*|j^*}, i^*$ (que on \unsure{a|n'a} précisément pas pour $f^*, i_*$) : il est bien vrai que si $X,Y$ sont des multistrates, $X',Y'$ leurs images, on a
\[
X\overset{\circ}{\ll}Y\Longrightarrow X'\overset{\circ}{\ll}Y'
\]
mais \uuline{pas}
\[
X\leq Y\Longrightarrow X'\leq Y'\quad\text{ou seulement }X'\ll Y'
\]
puisqu'il peut arriver que $Y'=\unsure{\varnothing|\phi},\ X'\ne\varnothing$ !

Cela montre donc que les \deleted{images}\inserted{applications} inverses $f^*$, pour $f$ \uline{non surjectif} \inserted{(celles \unsure{donc} qui ne sont pas des ``plongements'')}, \deleted{doivent être} \deleted{\unsure{traitées|traités} avec} mais qui impliquent une part des \unsure{passages} en question, doivent être traités avec \inserted{grande} attention. Ils ne sont pas justiciables d'une notion \unsure{pousse-bouton|presse-bouton} de ``morphisme'' de magasins ou d'atlas, où on aurait
\[
\begin{aligned}
X\overset{\circ}{\ll}Y&\Longrightarrow X'\overset{\circ}{\ll}Y'&&\text{mais aussi}\\
X\leq Y&\Longrightarrow X'\leq Y'&&\text{donc aussi}\\
X\ll Y&\Longrightarrow X'\ll Y'
\end{aligned}
\]
Mais dans tous les cas envisagés ici, \unsure{pourtant}
\[
X\mathrel{|\circ|}Y\Longrightarrow X'\mathrel{|\circ|}Y'.
\]
\scanpage{59}
\mspage{50}
% archive sheet number: 59
Le point délicat ici\unsure{,|;} est apparemment dans la notion de passage en \unsure{quotient}, \deleted{\unsure{plus cartes} ou contraction} qui nous fait sortir de la notion de ``homomorphisme \unsure{pousse-bouton|presse-bouton}.

Dans tous ces exemples, une multistrate est \unsure{transformée|transformé} en \unsure{item} (cas de \inserted{$f^*$ pour} $f$ surjectif, ou de $i_*$), ou bien en item ou $\unsure{\varnothing|\phi}$ (cas de $f^*$ en général). Il y a des cas qui s'apparentent au ``plongement'' d'un atlas dans un autre, où multistrate n'est \inserted{pourtant} pas \unsure{transformé|transformée} en item :

\uline{Exemple 2} Soit $L,L'$ deux ens, $L'\ne\varnothing$ considérons $\unsure{F''|F'}\in\mathrm{Figures}(L')$, \inserted{$F'\ne\varnothing$,} et
\diagram{\begin{tikzcd}[ampersand replacement=\&,column sep=large,row sep=large]
\mathrm{Fig}(L) \arrow[r,dash] \& \mathrm{Fig}(L\times L') \\
F \arrow[r,mapsto] \& F\times F'=\mathrm{pr}_1^*(F).\mathrm{pr}_2^*(F')
\end{tikzcd}}
C'est le composé d'un $\mathrm{pr}_1^*\deleted{\unsure{\circ}}$ (cas déjà regardé, $\mathrm{pr}_1$ surjectif, et d'une application interne à $\mathrm{Fig}(L\times L')$, savoir la
\[
\Phi\longmapsto\Phi.F\qquad\text{(où }F=\mathrm{pr}_2^*(F')\text{)}.
\]
C'est là un autre exemple d'une application qui (si $\operatorname{supp}F$ est $\unsure{\ne V|\ne\varnothing}$) a tendance à être un vrai plongement \inserted{\unsure{sauf|sans} \unsure{pour} figures (multistrates) \unsure{vides}}
\scanpage{60}
\mspage{51}
% archive sheet number: 60
\uline{Exemple 3} Passage d'un atlas \inserted{$\unsure{M|\mathcal{M}}$} à un sous-atlas \inserted{$M'$}, ou d'un magasin à un sous-magasin, dans l'esprit de restreindre le type de multistrates envisagées (par des propriétés du type connexité ou \unsure{inconnexité}, \uline{régularité}, être \deleted{\unsure{de}} des ``boules'' ou des ``simplexes'' dans un sens plus ou moins fort\unsure{''}) mais qui restent assez nombreuses pour ``\inserted{pouvoir} se \unsure{réduire} à $M'$'' -- notamment avec des \inserted{$M'$-}subdivisions \unsure{cofinales}.

Dans un tel cas, sûrement on demandera que pour le plongement $M'\longrightarrow M$, toutes les structures
\[
\leq,\quad\overset{\circ}{\ll},\quad\mathrel{|\circ|}\quad/\!/\quad\unsure{\gtrless|\bowtie},\quad\parallel
\]
soient conservées, à l'exception \unsure{tout} au plus de $\overset{\circ}{\ll}$ qui serait remplacé par $\underset{\mathrm{pol}}{\overset{\circ}{\ll}}$, et par conséquent $\gtrless$ par $\underset{\inserted{\mathrm{pol}}}{\gtrless}\deleted{\unsure{\mathrm{pol}}}$.

Donc on est ici dans le cas d'une notion d'hom \uline{de magasins}, qui est un plongement -- donc pas de difficultés du type de celles évoquées avec les \unsure{exemples|exemple} 1.
\scanpage{61}
\mspage{52}
% archive sheet number: 61
Cela nous confronte au passage avec la question de la stabilité des axiomes \inserted{des magasins}, par passage à un sous-magasin (\unsure{d'en entraînant|d'en impliquant}, pour un atlas associé au magasin, un sous-atlas correspondant).

Soit donc
\[
(\unsure{M|\mathcal{M}},\leq,\overset{\circ}{\ll},\mathrel{|\circ|})\quad\deleted{\text{un}}
\]
un magasin, $M'$ un sous-ens, que je munis des relations induites.
Je suppose
\[
\text{a)}\quad M'\ \text{fermé dans }M,\leq
\]
c'est la moindre des choses ! $\mathrm{Mg}\ \unsure{1^*|1^\circ}$ \inserted{(p. 22)} trivial, $\mathrm{Mg}\ 2^*$ a) b) aussi \inserted{grâce à a)}. Mais on voudrait que $\underset{M'}{\ll}$ \unsure{ce} soit $\underset{M}{\ll}\mathrel{|}M'$, ce qui amène à poser
\[
\text{b)}\quad\text{Si }X\overset{\circ}{\ll}Y'\leq Y,\quad\text{alors }X,Y\in M'\Rightarrow Y'\in M'
\]
$\mathrm{Mg}\ 3$ trivial, $\mathrm{Mg}\ 4$ aussi. Notons aussi que
\[
\parallel_{M'}=\parallel_M\mathrel{|}M,\qquad\unsure{\gtrless|\bowtie}_{M'}=\gtrless_M\mathrel{|}M.
\]
Tout marche !
\scanpage{62}
\mspage{53}
% archive sheet number: 62
\uline{Définition}\inserted{-prop.} Soit $\unsure{M|\mathcal{M}}$ un magasin, $M'$ une \deleted{sous-magasin}\inserted{partie de $M$}. On dit que c'est un \uline{sous-magasin} \unsure{s'il|si ça} satisfait
\[
\begin{aligned}
\text{a)}\quad&M'\ \text{fermé pour }\leq\qquad\text{i.e. }X\in M',\ Y\leq\deleted{\unsure{M|Y}}X\Rightarrow Y\in M'\\
\text{b)}\quad&\text{Si }X\overset{\circ}{\ll}Y'\leq Y\quad\text{et }X,Y\in M',\ \text{alors }Y'\in M'.
\end{aligned}
\]
Alors $M'$, muni des relations $\leq,\overset{\circ}{\ll},\mathrel{|\circ|}$ induites par $M$, est un magasin, et les relations $\ll,\unsure{\gtrless|\bowtie},\parallel$ \unsure{déduites} sont induites par les relations de $=$ nom dans $M$.

\uline{Définition} Dans ces conditions, si $\unsure{\mathscr{F}|\mathcal{F}|F}$ est un atlas associé à $M$, on appelle \uline{sous-atlas de $\mathscr{F}$ associé à $M'$} \deleted{\unsure{et}}, l'ens \inserted{$\mathscr{F}'$} \deleted{des $F\in\mathscr{F}$ tels que $\widetilde{F}\in$}
\[
\mathscr{F}'=\{F\in\mathscr{F}\mid\widetilde{F}\subset M'\}
\]
Pour la relation d'incidence induite \inserted{entre $M',\mathscr{F}'$} et les relations $\overset{\circ}{\ll},\mathrel{|\circ|}$ sur $M'$, c'est bien un atlas associé à $\deleted{\unsure{\mathscr{F}'|\mathscr{F}}}M'$. On a alors :

Les relations $\leq,\ll,\deleted{\unsure{\mathrel{|\circ|}|\parallel}}\gtrless,\parallel$ dans $\mathscr{F}'$ sont induites par les relations similaires de $\mathscr{F}$, le plongement : $\mathscr{F}'\longrightarrow\mathscr{F}$ commutant aux $\mathrm{Sup}$ quelconques pour $\leq$, et aux $\mathrm{Inf}$ des familles non vides.
\scanpage{63}
\mspage{54}
% archive sheet number: 63
Si $\mathrel{|\circ|}$ dans $\unsure{M|\mathcal{M}}$ peut être défini par $\overset{\circ}{\ll}$, il n'est pas clair qu'il en soit de $=$ dans $M$, mais ce sera le cas si $M'$ est $\overset{\circ}{\ll}$-fermé (mais \inserted{bien} souvent, $M'$ ne le sera pas\ldots). Par contre, souvent on aura $M'\supset\unsure{L|\mathcal{L}}$, de sorte que $L$ est aussi l'un des lieux de $M'$. Si alors $M$ est \unsure{str. dis|str. disc}, $M'$ aussi (axiomatique $\mathrm{Mg}$ \unsure{disc|dis}, page 23) ; \deleted{\unsure{et}} si $M$ est \unsure{quasi-ensembliste}, $M'$ aussi. Enfin si $M$ \unsure{ensembliste}, $M'$ aussi\ldots

\figure{\begin{tikzpicture}[x=1cm,y=1cm]
\draw[line width=.5pt] (0,0) .. controls (.18,-.16) and (.4,-.48) .. (.53,-.8);
\end{tikzpicture}}

Je suis mûr maintenant pour étudier la notion d'homomorphismes \inserted{(\unsure{restreints})} entre atlas, donc des hom. qui respectent, non seulement $\overset{\circ}{\ll}$ et $\mathrel{|\circ|}$, mais aussi $\leq$. On les appellera provisoirement ``homomorphismes''.

\deleted{\uline{Définition} Soit}

Soient donc
\[
\begin{aligned}
&(\unsure{\mathscr{F}|\mathcal{F}|F},M,\triangleleft,\overset{\circ}{\ll},\mathrel{|\circ|})\\
&(\mathscr{F}',M',\triangleleft,\overset{\circ}{\ll},\mathrel{|\circ|})
\end{aligned}
\]
deux atlas, on veut définir les applications
\[
f:\mathscr{F}\longrightarrow\mathscr{F}'
\]
qualifiées d'\uline{homomorphismes}.
\scanpage{64}
\mspage{55}
% archive sheet number: 64
$1^\circ)$ On veut qu'elle soit croissante pour $\leq$, et commute aux $\mathrm{Sup}$ quelconques. Ça implique qu'elle sera déterminée par sa restriction aux fig él.
\[
\unsure{f_M|f_\mu}:\unsure{M|\mathcal{M}}\longrightarrow\unsure{\mathscr{F}'|\mathcal{F}'|F'},
\]
qui sera une application croissante pour $\leq$. Si les figures de $\mathscr{F}$ sont de type fini, ou celles de $\mathscr{F}'$ \unsure{inconditionnelles}, n'importe quelle appl. \inserted{$f_M$} croissante pour $\leq$ proviendra d'un $f$.

$2^\circ)$ On veut qu'elle commute aux $\mathrm{Inf}$ quelconques, à familles \inserted{d'indices} non vides, et pour commencer
\[
f(F\cap G)=f(F)\cap f(G).\tag*{(*)}
\]
On est ramené aussitôt au cas où $F,G$ sont des multistrates, et la formule devient
\[
f_M(X)\cap f_M(Y)=\mathop{\mathrm{Sup}}_{Z\in\widetilde{X}\cap\widetilde{Y}}f_M(Z)\tag*{$(*_M)$}
\]
(qui rappelle la formule familière pour une figure \unsure{ensembliste}
\[
\unsure{X_i|x_i}\cap X_j=\bigcup_{k\leq i,j}X_{\unsure{n|k}}
\]
).

Cette formule (*) assure la commutation aux $\mathrm{Inf}$ finis. Pour les $\mathrm{Inf}$ infinis, on est ramené au cas des familles \deleted{\unsure{crois}} filtrantes décroissantes. Il n'est
\scanpage{65}
\mspage{56}
% archive sheet number: 65
pas clair que ce soit essentiel de le poser dans les définitions.

On veut d'autre part que $\unsure{f_M|f_\mu}$ soit compatible, en plus de $\leq$, à $\overset{\circ}{\ll},\mathrel{|\circ|}$. (\uline{NB} On sait déjà qu'il est compatible à $\unsure{\gtrless|\bowtie}$) Mais il faut lui donner un sens, dans le cas où \deleted{\unsure{on}} $f_M$ n'applique \inserted{pas} $\unsure{M|\mathcal{M}}$ dans $M'$, mais dans $\unsure{\mathscr{F}|\mathcal{F}|F}$ ! Mais il faut peut-être, avant cela, poser une autre condition, (comparer p. \unsure{24|28})

% vertical emphasis stroke beside the lower two lines of the condition
\[
\begin{aligned}
&\forall F\in\mathscr{F},\ \text{posant }F'=f(F),\ \text{l'application}\\
&\left|\begin{aligned}&G\longmapsto G'=f(F):\mathrm{Ss\,fig}(F)\overset{\sim}{\longrightarrow}\mathrm{Ss\,fig}\,F'\\
&\text{est un isom. d'ens ordonnés.}
\end{aligned}\right.\end{aligned}\tag*{(**)}
\]

\deleted{Utilisant le critère du \unsure{recollement|prolongement} des sous-figures, cela sera équivalent à :}
\[
\deleted{(**_M)\quad\forall X\in}
\]
Mais cela \unsure{indiquerait|impliquerait} que si $F$ est élémentaire, $F'$ aussi, c'est à dire finalement $f_M:M\longrightarrow M'$ ! C'est donc le moment de m'interroger si je connais des \deleted{\unsure{hom.}}\inserted{applications} $f:\mathscr{F}\longrightarrow\mathscr{F}'$ que j'aurais envie d'appeler ``homomorphismes'', et qui n'appliquent pas $M$ dans $M'$. Oui, il y a exemple \unsure{2|3} p. 50, mais du coup il m'apparaît hétéroclite etc.
\margin{\unsure{pas de quoi me laisser détourner de} (**).}
% margin note at the bottom edge, written horizontally
